\chapter{Elliptische Kurven}
    Das Kapitel basiert auf \cite{washington08} S.9ff.
    
    Beginnen wir also mit der Definition einer Elliptische Kurve.
    \begin{definition}
        Eine elliptische Kurve über einem Körper K wird definiert mit 
        \[
        y^2+a_1xy+a_3y=x^3+a_2x^2+a_4x+a_6,
        \]w
        wobei wir das die Weierstraß-Gleichung nennen.
        Wenn nun gilt char($K\neq 2,3$), so können wir die kurze
        Weierstraß-Gleichung nutzen
        \[
        y^2=x^3+Ax+B.
        \]
    \end{definition}
    Meistens wird eine \textbf{Elliptische Kurve} definiert mit: \[y^2=x^3+Ax+B,\]
    mit A und B Konstanten, dies bezeichnen wir als kurze Weierstraß-Gleichung für eine Elliptische Kurve. \(A,B,x\) und $y$ gehören dann zu einem Körper wie den rellen Zahlen \(\mathbb{R} \) oder auch den komplexen Zahlen $\mathbb{C}$. Es muss sich hierbei einfach im allgemeinen um einen Körper K handeln.
    Wir beschreiben nun die Menge der Punkte auf $L\subset K$ mit
    \[
    E(L)=\{\infty \} \cup \{(x,y)\in L\times L\mid y^2=x^3+Ax+B\},
    \]
    nach Definition ist der Punkt $\infty$ immer enthalten, warum wir diesen benötigen wird im nächsten Kapitel genauer erklärt. 
    Nicht über jedem Körper lassen sich Elliptische Kurven sinnvoll optisch darstellen, daher betrachten wir zuerst den Körper der reellen Zahlen. Hierbei gibt es zwei grundlegende Formen.
\begin{figure}
    \centering
    \begin{subfigure}{0.45\linewidth}
        \centering
        \includegraphics[width=\linewidth]{kap2_bs1.png}
        \caption{Eigene Darstellung mit GeoGebra $y^2=x^3-3x$}
    \end{subfigure}
    \hfill
    \begin{subfigure}{0.45\linewidth}
        \centering
        \includegraphics[width=\linewidth]{kap2_bs2.png}
        \caption{Eigene Darstellung mit GeoGebra $y^2=x^3+3x$}
    \end{subfigure}
    \label{fig:vergleich}
    \end{figure}
    Über dem Körper \(\mathbb{R}\) kann die Kurve je nach Anzahl der reellen Nullstellen verschiedene Formen annehmemen. Zum Beispiel hat $y^2 = x^3 - 3x$ drei verschiedene reelle Nullstellen, nämlich $x = 0$ und $x = \pm \sqrt{3}$. Im gegensatz dazu hat die Kurve $y^2 = x^3 + 3x$ lediglich eine reelle Nullstelle bei $x = 0$.\\

    Wie wir in der Darstellung sehen können, ändert sich die Form der Kurve über dem Körper $\mathbb{R}$, wenn wir mehrere reelle Nullstellen haben. Das Polynom \(x^3+Ax+B\) darf daher keine mehrfachen Nullstellen besitzen, dies ist äquivalent dazu, dass die Kurve \emph{nicht singulär} ist.

    \begin{definition}[(Einführung in die algebraische Gemoetrie, S.1+6)]
        Sei $f \in K[x,y]$ und
        \[
        E = V(f) = \{(x,y)\in K^2 \mid f(x,y)=0\}
        \] 
        eine affine ebene Kurve. Ein Punkt \((a,b)\in E\) heißt \textbf{singulär}, wenn,
        \[
        f_x(a,b)=0 \quad \text{und}\quad f_y(a,b)=0.
        \]
         Andernfalls heißt der Punkt \textbf{regulär}.

         Die Kurve \(E\) heißt \textbf{singulär}, wenn sie einen singulären Punkt besitzt.
    \end{definition}
    
    Wir schreiben die Gleichung nun um zu
    \[
    f(x,y)=y^2-x^3-Ax-B
    \]
    dann sind die Ableitung 
    \[
    f_x= -3x^2-A \quad \text{ und } \quad f_y=2y
    \]
    Es muss nun für beide ungleich 0 gelten,
    \[
    f_x= -3x^2-A = 0 \Leftrightarrow 3x^2+A = 0
    \]
    und
    \[
    f_y=2y = 0 \Leftrightarrow y = 0, \quad \text{char }\neq 2
    \]
    daraus folgt ein singulärer Punkt muss die Form $(x,0)$ haben und zusätzlich muss \(x^3+Ax+B = 0\) gelten, sodass der Punkt auf der Kurve liegt. Also x darf nicht eine Nullstelle von \(x^3+Ax+B\) und der Ableitung \(3x^2+A\) sein. Die Diskriminante für ein das kubische Polynom 
    \[
    p(x))x^3+Ax+B
    \]
    ist hierbei
    \[
    \Delta  = -4A^3-27B^2
    \]
   
    \[
        4A^3+27B^2 \neq 0.
    \]
    Für die Nullstellen $r_1, r_2, r_3$ des Polynoms gilt
    \[
    ((r_1 - r_2)(r_1 - r_3)(r_2 - r_3))^2 = -(4A^3 + 27B^2).
    \]
    Insbesondere folgt daraus, dass die Nullstellen genau dann paarweise verschieden sind, wenn die Diskriminante $\Delta = -(4A^3 + 27B^2)$ ungleich Null ist. Im vorliegenden Fall ist $\Delta < 0$, sodass voneinader verschiedene reelle Nullstellen vorliegen.
    
     
    \textbf{Allgemeine Weierstrass Gleichung}
    
    \[y^2+a_1xy+a_3y=x^3+a_2x^2+a_4x+a_6\]


\section{Gruppenstruktur}
Kapitel basier auf \cite{washington08}, S.12ff.
Basiert auch auf (An Introduction to the
Theory of Elliptic Curves, Silverman, S.8-19)
    Wir betrachten nun die Addition von zwei Punkten auf einer Elliptischen Kurve
   
    Wir betrachten dafür zwei Punkte $P=(x_1,y_1)$ und $Q=(x_2,y_2)$ auf einer elliptischen Kurve $E$. Diese ist wie bereits eingeführt gegeben durch die Gleichung \[y^2=x^3+Ax+B.\] Wir wollen nun einen dritten Punkt $P+_EQ=R $ berechnen. Für die graphische Darstellung wurde ein Beispiel über $\mathbb{R}$ verwendet mit der Kurve $y^2=x^3-4x+10$.\\

    Als ersten Schritt zeichnen wir eine Gerade $L$ durch $P$ und $Q$.
    \begin{figure}[htbp]
        \centering
        \includegraphics[width=0.5\linewidth]{kap2_2_bsp1.png}
        \caption{Gerade $L$ durch  $P$ und $Q$}
        \label{fig:placeholder}
    \end{figure}

    Als zweiten Schritt nehmen wir den Schnittpunkt $R'$ von der Gerade $L$ und $E$. Anschließend spiegeln wir $R'$ an der x-Achse und erhalten dadurch $R$. Wir definieren $R=P+_E Q$. Hierbei ist wichtig hervorzuheben, dass es sich dabei von der Addition der Koordinaten zweier Punkte unterscheidet.

    \begin{figure}[htbp]
        \centering
        \includegraphics[width=0.5\linewidth]{kap2_2_bsp2.png}
        \caption{Schnittpunkt $R'$ von Gerade $L$ und Kurve $E$ und die Spiegelung $R$}
        \label{fig:placeholder}
    \end{figure}

    Nun betrachten müssen wir noch den Fall betrachten, wenn wir einen Punkt $P$ auf sich selbst addieren wollen auf der Kurve $E$. Da es keinen zweiten Punkt, sodass durch beide Punkte eine Gerade gezogen werden kann, nehmen wir an der zweite Punkt $Q$ nähert sich P an, sodass die Gerade $L$ zu einer Tangente von $E$ bei $P$ wird.
    \begin{figure}[htbp]
        \centering
        \includegraphics[width=0.5\linewidth]{kap2_2_bsp3.png}
        \caption{Tangente am Punkt $P$}
        \label{fig:placeholder}
    \end{figure}
    
    Anschließend nehmen wir wieder den Schnittpunkt $R'$ von der Kurve $E$ und der Tangente $L$ und spiegeln diesen an der x-Achse. Dadurch erhalten wir wieder den Punkt $P+_EP=R$, denn wir aber nun als $2P$ bezeichnen.

    \begin{figure}[htbp]
        \centering
        \includegraphics[width=0.6\linewidth]{kap2_2_bsp4.png}
        \caption{piegelung von $R'$, wodurch $R=2P$ entsteht}
        \label{fig:placeholder}
    \end{figure}
    
    Als letzten Fall schauen wir uns noch an, wenn die Gerade $L$ keinen Schnittpunkt mit $E$ hat. Hierfür betrachten wir wieder den Punkt $P$ und seinen an der x-Achse gespiegelten Punkt $P'=-P$.
    \begin{figure}[htbp]
        \centering
        \includegraphics[width=0.5\linewidth]{kap2_2_bsp5.png}
        \caption{Schnittpunk von L im Punkt $\infty$}
        \label{fig:placeholder}
    \end{figure}
    Dafür brauchen wir einen dritten Punkte um $P+_E(-P)$ definieren zu kennen wir kreieren dafür einen Punkt $\infty$. Dieser Punkt ist dabei auf jeder vertikalen Gerade.\\

    \textbf{Herleitung der Additionsformeln
    \cite{washington08}(, S. 14 und Kaplan S.179ff.):}\\
    (Bei diesem Absatz habe ich mich orientiert an 
Kryptographie, Buell, S.99)
       
     Seien nun $P=(x_1,y_1)$ und $Q=(x_2,y_2)$ Punkte auf $E$ mit $P,Q\neq \infty$.
    Zuerst betrachten wir den Fall $P\neq Q$ und \(x_1\neq x_2\). Dann hat die Gerade \(L\) durch durch \(P\) und \(Q\) die Form
    \[
    y=m(x-x_1)+y_1
    \qquad \text{mit} \qquad
    m=\frac{y_2-y_1}{x_2-x_1}.
    \]
    Zur Bestimmung des dritten Schnittpunkts setzen wir die Geradengleichung in die Kurvengleichung
    \[
    E: y^2=x^3+Ax+B
    \]
    ein. Es gilt also
    \[
    \bigl(m(x-x_1)+y_1\bigr)^2=x^3+Ax+B.
    \]
    Nach Ausmultiplizieren und Umstellen erhält man
    \[
    0=x^3-m^2x^2+\bigl(2m^2x_1-2my_1+A\bigr)x
    +\bigl(B-m^2x_1^2+2mx_1y_1-y_1^2\bigr).
    \]
    
   Die Nullstellen des kubischen Polynoms sind die \(x\)-Koordinaten der drei Schnittpunkte von \(L\) mit \(E\), also \(x_1\), \(x_2\) und \(x_3'\). Daher gilt
    \[
    0=(x-x_1)(x-x_2)(x-x_3')
    =x^3-(x_1+x_2+x_3')x^2+(x_1x_2+x_1x_3+x_2x_3)x-x_1x_2x_3
    \]
    Mit Kooeffizientenvergleich erhält man
    \[
    x_1+x_2+x_3' = m^2 \Leftrightarrow  x_3' = m^2-x_1-x_2
    \]
    Da \(R'=(x_3',y_3')\) auf der Geraden \(L\) liegt, gilt
    \[
    y_3'=m(x_3'-x_1)+y_1.
    \]
    Die Summe \(P+Q\) ist die Spiegelung von \(R'\) an der \(x\)-Achse. Somit ist
    \[
    P+Q=(x_3,y_3)
    \]
    mit
    \[
    x_3=x_3'=m^2-x_1-x_2,
    \qquad
    y_3=-y_3'=m(x_1-x_3)-y_1.
    \]    
    
    \textbf{voll restliche Herleitung Einfügen, Rosen S.12-14}
    
    Dann definieren wir $P+Q=R=(x_3,y_3)$ wie folgt:
    \begin{enumerate}
        \item Wenn $x_1\neq x_2$, dann
        \[x_3=m^2-x_1-x_2, \quad y_3=m(x_1-x_3)-y_1 \quad \text{ mit  } m=\frac{y_2-y_1}{x_2-x_1}.\]
        \item Wenn $x_1=x_2$, aber $y_1\neq y_2$, dann ist 
        \[P+Q=\infty.\]
        \item Wenn $P=Q$ und $y_1\neq 0$, dann
        \[x_3=m^2-2x_1, \quad y_3=m(x_1-x_3)-y_1, \quad \text{ mit } m=\frac{3x_1^2+A}{2y_1}.\]
        \item Wenn $P=Q$ und $y_1=0$ dann, 
        \[P+Q=\infty.\]
    \end{enumerate}

    Zusätzlich definieren wir 
    \[P+\infty=P\]
    für alle Punkte $P$ auf $E$.Damit ist $E(L)$ geschlossen unter Addition von Punkte auf der Kurve. 

    \begin{theorem}
        Die Addition auf einer Elliptische Kurve $E$ erfüllt folgende Eigenschaften.
        \begin{enumerate}
            \item [Existenz der Identität:] \(P+\infty=\infty+P=P\) für alle \(P\in E\)
            \item [Kommutativität:] \(P+Q=Q+P\) für alle \(P,Q\in E\)
            \item [Existenz des Inversen:] Sei $P\in E$ gegeben, dann exestiert ein $-P\in E$ mit \(P+(-P)= \infty\), für alle $P\in E$.
            \item [Assoziativität:] \((P+Q)+R=P+(Q+R)\) für alle \(P,Q,R\in E\)
        \end{enumerate}
    \end{theorem}
    Damit formen alle Punkte auf $E$ eine additive abelsche Gruppe mit $\infty$ als neutrales Element.

\begin{proof}
    Die Existenz der Identität ist durch die Definition der Elliptische Kurve gegeben. Die Kommutativität kann einfach nachgerechnet werden indem die oben definieren Regeln für die Addition zweier Punkte anwenden.
    Für  das Inverse müssen wir von einem Punkt $P$ sein Reflektion an der x-Achse $P'$ nehmen. Dann erhalten wir $P+P'=\infty.$
    Der Beweis der Assoziativität ist recht langwierig und daher wird auf eine externe Quellen hierfür verwiesen \textbf{Quelle für Assoziativität einfügen, Potenziell (\cite{washington08},S.20ff. Proof of Associativity}
\end{proof}


\textbf{Beispiel 1:}
    Wir betrachten dafür eine Elliptische Kurve $E$ mit 
    \[y^2=x^3-4x+4,\]
    auf dieser Kurve liegen zum Beispiel die Punkte $P=(0,2)$ und $Q=(2,2)$. Nun berechnen wir auf den vorherigen Regeln, da $x_1\neq x_2$ folgt,
    \[
        m = \frac{2-2}{2-0}=0
    \]
    \[
        x_3 = (0)^2-0-2=-2, \quad  y_3 = 0\cdot(0-(-2))-2 =-2
    \]
    also
    \[
    P+_EQ=(0,0)+_E(0,2)= (-2,-2)
    \]
    Analog können wir beliebige Punkte auf der Kurve addieren.
    \[(0,2)+_E(-2,-2)=(6,14)\]

    
    \textbf{Beispiel 2:}
    Nun können wir auch Punkte berechnen, wenn wir diese geometrisch nicht sinnvoll darstellen können. Die Kurve \(y^2=x^3-4x+4\pmod 7\) also über dem endlichen Körper $\mathbb{F}_7$. (Als Errinnerung die Kurve muss über dem Körper Singulär sein\(4A^3+27B^2=4\cdot(-4)^3+27\cdot4^2=176\equiv1\pmod 7\neq0\)

    Wir wählen nun die Punkte $P=(0,2)$ und \(Q=(5,5)\)
    auf $E(\mathbb{F})_7$.
    Es gilt nun 
    \[
    m=\frac{5-2}{5-0}=\frac{3}{5} = 3\cdot 5^{-1}\equiv 3\cdot 3 \pmod 7= 9 = 2 \pmod 7
    \]
    und damit
    \[
    x_3=2^2-0-5=-1\equiv 6 \pmod 7, \quad y_3=2(0-6)-2=-14 \equiv 0\pmod 7
    \]
    also
    \[P+_EQ=(6,0).\]

    
    \textbf{Skalar mal Punkt}(The aglebra of Elliptic Curves, S.30f.)
    Um später Diffie-Hellman über der endlcihen Gruppe $E(\mathbb{F}_p)$ anwenden zu können(wird im Kryptographie Kapitel genauer erklärt), muss man auch für größere $p$ berechnen können, also
    \[mP=P+P+\dots+P\in E(\mathbb{F}_p)\] für sehr große Werte von m.
    Lass \(m\in\mathbb{N}_{>0}\).

    Wir schreiben zuerst 
    \[m=m_0+m_1\cdot2+m_2\cdot 2^2+\dots+m_r\cdot 2^r, \quad m_0,\dots,m_r\in\{0,1\}\]
    also zerlegen m in seinen einzelnen Zweierpotenzen. Dann kann $mP$ berechnet werden als
    \[mP=m_0P+m_1\cdot 2P+m_2\cdot2^2P+\dots +m_r\cdot 2^rP,\]
    wobei \(2^kP=2\cdot 2\cdots 2P\) benötigt nur k Doppelungen und damit können wir $mP$ in $O(\text{log } m)$ berechnen. Die Binärdarstellung von $m$ hat 
    \[
    r+1 = \lfloor\log_2m \rfloor+1
    \].
    Im Schnitt braucht man also $\log_2(m)$ Verdoppelungen und \(\frac{1}{2}\log_2(m)\) Additionen um $mP$ zu berechnen.

    Man kann die benötigte Zeit nochmal reduzieren, indem man die Non-Adjacent Form verwendet. Die Addition und Subtraktion eines Punktes kosten genau gleich viel Zeit, also schreiben wir
    \[
    m=m_0+m_1\cdot 2+m_2\cdot 2^2+\cdots+m_r\cdot 2^r, \quad m_0,\dots,m_r\in\{-1,0,1\}.
    \]
    Hierfür können wir im Schnitt  $\frac{2}{3}$ der $m_i=0$ für die Zerlegung nutzen, dass reduziert die Anzahl der durchschnittlichen Additionen auf \(frac{1}{3}\log_2(m)
    \).
    Aufgrund dieses Vorgehens ist die Berechnung von $m$ aus $P$ und $mP$ sehr schwierig. Das sogennante \textbf{diskrete Logarithmus Problem} für Elliptische Kurven ist später essenziell für das kryptographische Verfahren, das auf diesen aufbaut.
   

\section{Projektiver Raum und der Punkt bei Unendlich} (vgl. \cite{washington08}, S. 18ff.)
(vgl. \cite{clader25}, S.245f.)
    
    Die Gruppenstruktur elliptischer Kurven wurde zuvor über die geometrische
    Punktaddition beschrieben. Dabei spielt der Punkt im Unendlichen eine besondere
    Rolle, da er als neutrales Element der Gruppe dient. Um diesen Punkt geometrisch
    präzise zu beschreiben, gehen wir nun von der affinen Ebene zur projektiven Ebene über. 

    \begin{definition}[Projektive Ebene]
    Sei \(K\) ein Körper. Die projektive Ebene über \(K\) ist definiert als der Quotient
    \[
    \mathbb P^2_K
    =
    \left(K^3\setminus0\right)/\sim,
    \]
    wobei zwei Tripel \((x,y,z),(x',y',z')\in K^3\setminus\{(0,0,0)\}\) genau dann in einer Relation  stehen \(
    (x,y,z)\sim (x',y',z)\), wenn ein \(\lambda\in K^\times:
    (x,y,z)=\lambda(x',y',z')\) exestiert.
    Die Äquivalenzklasse von \((x,y,z)\) wird mit
    \((x:y:z)\) bezeichnet und heißt Punkt der projektiven Ebene.
    \end{definition}
    Die Koordinaten \((x:y:z)\) werden homogene Koordinaten genannt. Dabei bezeichnet \((x:y:z)\) nicht einen einzelnen Punkte, sondern die Äquivalenzklasse aller von Null verschiedenen skalaren Vielfachen von \((x,y,z)\). Für jedes \(\lambda\in K^\times\) gilt also
    \[
    (x:y:z)=(\lambda x:\lambda y:\lambda z).
    \]
     Die Menge 
    \[
    g_\infty=\{(x:y:z)\in \mathbb P^2_K\mid z=0\}.
    \]
    heißt die unendlich ferne Gerade der projektiven Ebene (vgl. \cite{forster15}, S.200ff.).
    Ist $z \neq 0$, so kann z auf 1 normiert werden, denn es gilt \((x:y:z)=(x/z:y/z:1)\). Nun kann man eine Abbildung definieren
    \[
    K^2\longrightarrow\mathbb{P}^2_K, \quad
    (x,y)\mapsto (x:y:1).
    \]
    Diese Punkte entsprechen genau den affinen Punkten \((x,y)\in K^2\) die bijektiv auf \(\mathbb{P}^2_K\setminus g_\infty\) abgebildet werden.


    \begin{lemma}
        Die Punkte auf der Geraden im Unendlichen \(L_\infty\) entsprechen den RIchtungen affiner Geraden. Insbesondere schneiden sich parallel affine Geraden in der projektiven Ebene in demselben Punkt auf \(L_\infty\).
    \end{lemma}
    Ich weiß nicht ganz wie ich das definieren soll.

    Soll ich auch definieren, dass es sich um eine Äquivalenzklasse handelt?
\section{Beweis der Assoziativitität}
    Beweis nach vgl. \cite{zwegers24}, S. 477ff.
    Es gibt hierbei gibt es 3 Sonderfälle, erstens die Punkte könnten in unendlich liegen, dafür brauchen wir die projektiven Koordinaten. Zweitens die Gerade könnte eine Tangente bei E sein, was bedeutet, dass zwei Punkte gleich sind. Drittens zwei Geraden könnten gleich sein.
    Wir möchten nun die Assoziativität zeigen. Sei \(E=y^2=x^3+Ax+B\) eine elliptische Kurven über einem Körper \(K\) mit \(\text{char(K)}\neq 2,3\) und \(4A^3+27B^2 \neq 0\).
    Weiter seien nun die Punkte\[
    P_1=(x_1,y_1),\quad P_2=(x_2,y_2),\quad P_3=(x_3,y_3)
    \]
    Punkte aus \(E\setminus\{O\}\). Ziel ist es zu zeigen, dass
    \[
    (P_1+P_2)+P_3=P_1+(P_2+P_3)
    \]
    gilt.

   Für zwei Punkte \(P_i=(x_i,y_i)\) und \(P_j=(x_j,y_j)\) mit
    \(P_j\neq -P_i\) setzen wir
    \[
    \alpha(P_i,P_j):=
    \begin{cases}
    \dfrac{y_j-y_i}{x_j-x_i}, & \text{falls } x_i\neq x_j,\\[1.2ex]
    \dfrac{3x_i^2+A}{2y_i}, & \text{falls } P_i=P_j \text{ und } y_i\neq 0.
    \end{cases}
    \]
    Wir suchen nun eine Parabel durch \(P_1,P_2,P_3\). Wir finden dafür einen neuen Punkt \(P_4=(P_1+P_2)+P_3\), definiert mit
    \[
    x_4=-x_1+x_2-x_3+\frac{1}{c_2^2}-2\frac{c_1}{c_2}, \quad y_4 = -y_2-(x_4-x_2)(c_1+x_4c_2)
    \]
    wobei 
    \[
    c_1 := \frac{x_1\alpha(P_2,P_3)-x_3\alpha(P_1P_2)}{x_1-x_3}, \quad c_2:=\frac{\alpha(P_1,P_2)-\alpha(P_2,P_3)}{x_1-x_3}.
    \]

   
    

\section{Elliptische Kurven in Charakteristik 2} (nach Rosen, S. 47ff.) \textbf{Muss man hier auch characteristik 3 ergänzen?}

    Computer arbeiten tendenziell besser über Körpern der Charakteristik 2 im Vergleich zu Körpern mit großen prim Charakteristiken.Daher werden häufig Elliptische Kurven über einem Feld $\mathbb{F}_q$ mit $q=2^k$ Elementen verwendet.
    (vgl. \cite{silverman06}, S.47)

    Aus dieser Notwendigkeit betrachten wir nun genauer Elliptische Kurven in Charakteristik 2, da hier einige unserer vorher definierten Formeln nicht funktionieren.\\
    \begin{satz}
         In Charakteristik \(2\) ist die kurze Weierstrass Gleichung \(y^2=x^3+Ax+B\) stets singulär. Daher verwendet man in Charakteristik 2 die allgemeine Weierstraß-Gleichung
    \end{satz}

    \[
    y^2+a_1xy+a_3y=x^3+a_2x^2+a_4x+a_6.
    \]
    \begin{proof}    
        Betrachte 
        \[
        f(x,y) = y^2-x^3-Ax-B.
        \]
        Ein Punkt auf der Kurve ist genau dann singulär, wenn beide partiellen Ableitungen verschwinden. Es gilt
        \[
        f_y=2y\equiv0,
        \] da der Körper Charakteristik 2 hat. Außerdem ist 
        \[
        f_x=-3x^2-A
        \]
        Sei nun $x_0$ eine Nullstelle von $f_x$. Wir wählen ein \(y_0\) mit
        \[
        y_0^2=x_0^3+Ax_0+B
        \]
        Dann liegt $(x_0,y_0)$ auf der Kurve und es ist
        \[
        f_x(x_0,y_0)=f_y(x_0,y_0) = 0
        \]
        Damit folgt die Kurve ist singulär.
    \end{proof}
    
    Nun wollen auch für Charakterstik 2 die Gruppengesetze aufstellen.
    Hierfür beginnen wir mit der allgmeinen Weierstraß-Gleichung für eine elliptische Kurv $E$
    \[
    y^2+a_1xy+a_3y=x^3+a_2x^2+a_4x+a_6.
    \]
    Man schreibt diese in die Form
    \[
    f(x,y)=y^2+a_1xy+a_3y-x^3-a_2x^2-a_4x-a_6.
    \]
    Dann sind die partiellen Ableitungen gegeben durch
    \[
    f_x=a_1y-3x^2-2a_2x-a_4,\qquad f_y=2y+a_1x+a_3.
    \]
    In Charakteristik \(2\) vereinfacht sich dies zu
    \[
    f_x=a_1y+x^2+a_4,\qquad f_y=a_1x+a_3.
    \]
    Ein Punkt \((x,y)\) ist wie bereits früher eingeführt, genau dann singulär, wenn
    \[
    f(x,y)=0,\qquad f_x(x,y)=0,\qquad f_y(x,y)=0.
    \]
    Fall 1: $a_1\neq 0$
    Man kann aus \(a_1x+a_3=0\) folgern
    \[
    x=\frac{a_3}{a_1}.
    \]
    Eingesetzt in \(a_1y+x^2+a_4=0\) liefert es,
    \[
    a_1y+\left(\frac{a_3}{a_1}\right)^2+a_4=0
    \]
    \[
    y= \frac{a_3^2}{a_1^3}+ \frac{a_4}{a_1}.
    \]
    wenn der Punkt zusätzlich \(f(x,y)=0\) erfüllt handelt es sich um einen singulären Punkt.
  
    Fall 2: $a_1=0$
    Es gilt nun,
    \[
    f_x=x^2+a_4,\qquad f_y=a_3.
    \]
    also muss \(a_3= 0\) und $f(x,y)=0$ gelten, sodass $(x,y)$ ein singulärer Punkt ist.\\

    Nun können wir die Addition von Punkten definieren. Es gilt wieder für alle Punkte $P\in E$
    \[
    P+\infty=P.
    \]
    \textbf{washington S.48-49 hier weiter machen}

